% =====================================================================
\chapter{事後処理}
% =====================================================================

当日が終わっても計センの仕事は終わりません。
ここで何を残すかによって，次の担当者の負担が変わります。

\begin{Fig}
\begin{tikzpicture}[x=1mm,y=1mm,every node/.style={font=\sffamily\scriptsize}]
  \draw[line width=1pt,rulec] (0,0)--(150,0);
  \foreach \x/\t/\d in {
      6/{当日中}/{Step 38\\判定を確定},
     42/{当日〜翌日}/{Step 39\\ラップ公開},
     78/{1 週間以内}/{Step 40\\公式成績表\\Step 41 保存},
    114/{返却期日まで}/{Step 42\\機材返却\\Step 43 振り返り}}{
    \fill[boxln] (\x,0) circle (1.2);
    \node[font=\sffamily\bfseries\scriptsize,above=2mm of {(\x,0)},anchor=south] {\t};
    \node[font=\sffamily\tiny,below=2.5mm of {(\x,0)},anchor=north,align=center] {\d};
  }
  \node[anchor=north,font=\sffamily\tiny,align=center] at (75,-16)
    {※ Step 38 で判定が確定していないと，公開も成績表も作れません};
\end{tikzpicture}
\figcap{事後処理の流れと期限}
\end{Fig}

% ---------------------------------------------------------------------
\Step{成績を確定させる}{\tALL}

\begin{Goal}
すべての競技者の判定を確定し，「保留」が 1 件も無い状態にする。
\end{Goal}

\begin{enumerate}
\item メインウィンドウの検索窓に \textbf{DISQ}，\textbf{DNS}，\textbf{DNF}
      と入力して，該当者を一覧で確認します。
      \textbf{検索窓はナンバー・氏名の一部・所属・判定のいずれでも引けます。}
\item 未確定・保留のまま残っている競技者がいないかを確認します。
\item 判定に迷ったものは，\textbf{競技責任者に判断してもらいます}。
      計センは事実（何番が記録されているか，バックアップに何があるか）を提示する役です。
\item 全クラスのリザルトリストを表示し，順位が正しく付いていることを確認します。
      \textbf{競技時間を超えた人に順位が付いていないか}も見ます。
\end{enumerate}

\begin{Done}
全競技者の判定が確定し，リザルトリストの人数がエントリー数と一致している。
\end{Done}

% ---------------------------------------------------------------------
\Step{ラップと成績を公開する}{\tALL}

\begin{Goal}
参加者がラップを見られる状態にする。
\end{Goal}

\subsection*{(1) Lap Center へ掲載する}

\begin{enumerate}
\item メインウィンドウ →【ファイル出力】→\textbf{【Lap Center へアップロード】}を選ぶ。\ver
\item \textbf{【記録一覧／ラップ解析】}を選び，クラスを全選択して【出力】。
\item 次の画面で\textbf{【直接アップロード】}を選ぶ。
      \textbf{Mulka2 のライセンスが必要}です。
\item 掲載アカウントの情報を入力する。
\item \textbf{アップロード先}（一般公開のフォルダ／限定のフォルダ）を選び，【実行】。
\end{enumerate}

\begin{Note}[アップロードできないとき]
Mulka2 は\textbf{全角と半角の区別に弱い}ところがあります。
アカウント名に含まれる文字（小文字の l と大文字の I など）の\textbf{見分けがつきにくい文字}が
原因で認証に失敗することがあります。うまくいかないときは，
\textbf{アカウント名を打ち直す}，\textbf{パスワードを入れ直す}を試します。
\end{Note}

\subsection*{(2) ライブ速報（競技中に公開する場合）}

クラウドを使っている場合，【通信マネージャ】→【成績公開設定】から
成績のライブ速報を出せます。\ver

\begin{itemize}
\item \textbf{一般公開}：誰でも見られます。
\item \textbf{限定公開}：URL を知っている人だけが見られます。
\item \textbf{公開のタイミングは競技責任者と相談}して決めます。
\end{itemize}

\begin{Fail}
\textbf{起きること：}スタート閉鎖前にライブ速報を公開してしまう。\\
\textbf{対処：}ライブ速報には\textbf{未帰還者の表示}が含まれることがあり，
まだ出走していない競技者に\textbf{他人の成績が見えてしまいます}。
セレクションを兼ねた大会では特に問題になります。
\textbf{公開はスタート閉鎖後にする}のが安全です。
参加者を隔離できない大会では，実況・速報・ライブを行わない選択もあります。
\end{Fail}

\subsection*{(3) 内部限定で公開する場合}

一般公開したくない場合（部内戦など）は，
【ファイル出力】→【その他】からラップ解析用のファイルを出し，
HTML に変換して内部で共有する方法があります。\ver

% ---------------------------------------------------------------------
\Step{公式成績表を作る}{\tRANK}

\begin{Goal}
公認大会の要件を満たした公式成績表を作り，所定の提出先に出す。
\end{Goal}

\begin{enumerate}
\item メインウィンドウ →【ファイル出力】→【その他】→
      \textbf{【競技者結果データ一覧（CSV 形式）】}でデータを出す。\ver
\item 出力はゼッケン番号順なので，\textbf{クラス順・タイム順に並べ替えます}。
\item 次の項目を抽出します。
      \begin{itemize}
      \item 順位／記録（所要時間）／氏名／所属／\textbf{競技者登録番号}
      \end{itemize}
\item 大会名・開催日・テレイン名・クラス名・コース距離・登距離・コントロール数などの
      ヘッダを付けます。
\item \textbf{アドバイザ（大会役員）の確認を受けます}。
\item 所定の提出先に提出し，大会サイトにも掲載します。
\end{enumerate}

\begin{Note}[期限]
公式成績表は\textbf{大会後 1 週間以内}を目標にします。
記憶が新しいうちに作らないと，判定の根拠が分からなくなります。
\end{Note}

\begin{Done}
公式成績表ができ，確認を受けて提出した。
\end{Done}

% ---------------------------------------------------------------------
\Step{データを保存・アーカイブする}{\tALL}

\begin{Goal}
「あとから問い合わせが来ても答えられる」状態でデータを残す。
\end{Goal}

大会後に「自分の記録がおかしい」という問い合わせが来ることがあります。
そのときに\textbf{当日のデータがそのまま残っていること}が唯一の根拠になります。

\begin{center}
\small
\begin{tabular}{@{}P{44mm}P{100mm}@{}}
\toprule
\hd{保存するもの} & \hd{内容・注意} \\
\midrule
イベントデータフォルダ一式 & \texttt{Changes.dat} を含めて\textbf{フォルダごと}。
  これがあれば当日の全操作を再現できます \\
元のエントリーデータ & 変換前の生データと，変換設定ファイル \\
スタートリスト（全 4 種類） & 実際に配った版。版数入り \\
コースデータ & 当日に変更した場合は，変更後の内容もメモ \\
ピンパンチ対応表 & 写真でも可 \\
\tSI\ ステーションのバックアップメモリ & トラブルがあった場合はテキストで保存 \\
当日の申込用紙・変更届 & 現物をまとめて保管 \\
トラブルの記録 & 何が起きて，どう対応したか。次回の申し送りになります \\
\bottomrule
\end{tabular}
\end{center}

\begin{Done}
上記が 1 つのフォルダにまとまり，複数人の手元にコピーがある。
\end{Done}

% ---------------------------------------------------------------------
\Step{機材を整理して返却する}{\tALL\ \tSI\ \tEMIT\ \tSIAC}

\begin{Goal}
借りた状態に戻して返す。次に借りる人が困らないようにする。
\end{Goal}

\begin{enumerate}
\item \textbf{カード}：問題がなくなった時点でラベルをはがし，
      汚れていれば洗って乾かしてから返却します。
      \tSIAC\ 電源が入ったままの SIAC は，SIAC OFF ステーションで切ります。
\item \tSI\ \textbf{ステーション}：\textbf{Service/OFF カードでスタンバイモードに戻します}。
      とくに\textbf{タッチフリー設定のステーションは通常の約 10 倍の電力を消費する}ので，
      撤収したらすぐに戻します。
\item \tSI\ \textbf{設定を変更した機材は，原則として元の設定に戻します}。
      タッチフリー用に変更した機材を通常設定に戻す，
      番号を変更した機材を元の番号に戻す，など。
      デフォルト設定を持っている機材なら，デフォルトに戻す操作で一括して戻せます。
\item \textbf{パンチ台}：泥を洗い流し，拭いてから返却します。
\item \tSI\ 雨の大会のあとは，\textbf{ステーション内部への浸水}がないか確認します。
      封止不良で内部に水滴が入り，パンチ時刻が不正に記録された事例があります。
\item \textbf{枚数・台数を数えて}，貸出元に返却報告をします。
      使用枚数の報告が必要な場合は，備品担当に伝えます。
\end{enumerate}

\begin{Done}
借用リストの全項目が返却済みになり，返却報告をした。
\end{Done}

\begin{Fail}
\textbf{起きること：}番号を変更したステーションを，そのまま返却する。\\
\textbf{対処：}次の大会でその機材を使う人が，\textbf{ラベルと中身が違う}状態で使うことになります。
番号を変更したら\textbf{必ず記録し，返却前に戻します}。
戻せない場合は，\textbf{何番を何番に変えたかを書いたメモを添えて}返却します。
\end{Fail}

% ---------------------------------------------------------------------
\Step{振り返りを残す}{\tALL}

\begin{Goal}
次の大会の担当者が，同じ失敗を繰り返さないようにする。
\end{Goal}

大会後 1 週間以内に，次の 3 点を書き残します。長文にする必要はありません。

\begin{enumerate}
\item \textbf{起きたトラブルと，実際にやった対応}（うまくいったか・いかなかったかも）
\item \textbf{事前準備で足りなかったもの}（機材・時間・人手・打合せ）
\item \textbf{次回やる人へ渡す具体的な数字}
      \begin{itemize}
      \item 参加者数に対して必要だった読み取り台数・人数
      \item レンタルカードの使用枚数
      \item ステーションの動作時間の設定値と，実際に足りたかどうか
      \item 各作業に実際にかかった時間
      \end{itemize}
\end{enumerate}

\begin{Note}[パートマニュアルを更新する]
本書のようなテキストとは別に，\textbf{自分の大会用のパートマニュアル}を持っているはずです。
それを更新して次に渡すのが，最も効きます。
とくに\textbf{機材リストの数量}と\textbf{当日のタイムテーブル}は，
実績値に置き換えます。
\end{Note}

\begin{Done}
振り返り文書ができ，次期担当者に渡した。パートマニュアルを更新した。
\end{Done}
